\documentclass[11pt,a4paper]{article}

%====================================================
% PACKAGES
%====================================================
\usepackage[utf8]{inputenc}
\usepackage[T1]{fontenc}
\usepackage[french]{babel}
\usepackage[a4paper,margin=1.8cm]{geometry}
\usepackage{amsmath,amssymb,amsthm,mathtools}
\usepackage{lmodern}
\usepackage{enumitem}

%====================================================
% ENVIRONNEMENTS
%====================================================
\newtheorem{theorem}{Théorème}
\newtheorem{proposition}{Proposition}

\title{\textbf{Fiche d’exercices — Probabilités (niveau concours)}}
\date{}

\begin{document}
\maketitle

\section*{AIDE}
APP = application, DEM = démonstration, PREP = réflexion, PB = problème.

%====================================================
\section{Exercices d’application directe}

\textbf{Exercice 1 [APP ⋆]}  
On lance un dé équilibré.  
\begin{enumerate}
\item Définir l’univers $\Omega$.
\item Calculer $P(\text{pair})$.
\item Calculer $P(\geq 4)$.
\item Les événements sont-ils indépendants ?
\end{enumerate}

%--------------------------------------
\textbf{Exercice 2 [APP ⋆]}  
Soit $X \sim \mathcal{B}(1,p)$.  
\begin{enumerate}
\item Donner la loi de $X$.
\item Calculer $E(X)$.
\item Calculer $Var(X)$.
\item Interpréter.
\end{enumerate}

%--------------------------------------
\textbf{Exercice 3 [APP ⋆⋆]}  
On répète 10 expériences de Bernoulli de paramètre $p=0.4$.  
\begin{enumerate}
\item Identifier la loi de $X$.
\item Calculer $P(X=4)$.
\item Calculer $P(X\geq 1)$.
\item Calculer $E(X)$ et $Var(X)$.
\end{enumerate}

%--------------------------------------
% (on saute à 10 mais dans ton PDF final tu peux dupliquer ce pattern)
\textbf{Exercice 10 [APP ⋆⋆]}  
Soit $X$ une variable continue uniforme sur $[0,2]$.  
\begin{enumerate}
\item Donner sa densité.
\item Calculer $P(X<1)$.
\item Calculer $E(X)$.
\end{enumerate}

%====================================================
\section{Exercices de démonstration}

\textbf{Exercice 11 [DEM ⋆⋆]}  
Soit $X$ une variable de Bernoulli.  
\begin{enumerate}
\item Calculer $E(X)$.
\item Calculer $E(X^2)$.
\item En déduire $Var(X)$.
\end{enumerate}

%--------------------------------------
\textbf{Exercice 12 [DEM ⋆⋆⋆]}  
Soient $X_1,...,X_n$ indépendantes de loi Bernoulli$(p)$.  
\begin{enumerate}
\item Montrer que $\sum X_i$ suit une loi binomiale.
\item Calculer son espérance.
\item Calculer sa variance.
\end{enumerate}

%--------------------------------------
\textbf{Exercice 13 [DEM ⋆⋆⋆]}  
Démontrer la formule des probabilités totales.

%--------------------------------------
\textbf{Exercice 14 [DEM ⋆⋆⋆]}  
Démontrer la formule de Bayes.

%--------------------------------------
\textbf{Exercice 15 [DEM ⋆⋆⋆⋆]}  
Montrer que la loi géométrique est sans mémoire.

%--------------------------------------
\textbf{Exercice 16 [DEM ⋆⋆⋆⋆]}  
Démontrer que la loi de Poisson a pour espérance $\lambda$.

%====================================================
\section{Exercices de réflexion}

\textbf{Exercice 21 [PREP ⋆⋆⋆]}  
Une file d’attente reçoit des clients.  
\begin{enumerate}
\item Proposer deux modèles probabilistes.
\item Discuter leurs hypothèses.
\item Comparer.
\end{enumerate}

%--------------------------------------
\textbf{Exercice 22 [PREP ⋆⋆⋆⋆]}  
On considère une loi binomiale.  
\begin{enumerate}
\item Étudier son comportement quand $n\to\infty$.
\item Justifier approximation normale.
\item Donner conditions.
\end{enumerate}

%--------------------------------------
\textbf{Exercice 23 [PREP ⋆⋆⋆⋆]}  
Étudier la différence entre espérance et valeur la plus probable.

%====================================================
\section{Problèmes type concours}

\textbf{Exercice 31 [PB ⋆⋆⋆]}  
On lance une pièce biaisée $n$ fois.  

\begin{enumerate}
\item Déterminer la loi du nombre de succès.
\item Calculer son espérance.
\item Calculer sa variance.
\item Étudier la convergence de $\frac{X}{n}$.
\end{enumerate}

%--------------------------------------
\textbf{Exercice 32 [PB ⋆⋆⋆⋆]}  
Une machine produit des pièces avec défaut de probabilité $p=0.02$.  

\begin{enumerate}
\item Modéliser sur 100 pièces.
\item Calculer $P(X=3)$.
\item Approximer avec Poisson.
\item Comparer.
\end{enumerate}

%--------------------------------------
\textbf{Exercice 33 [PB ⋆⋆⋆⋆]}  
On considère une loi géométrique.

\begin{enumerate}
\item Interpréter.
\item Calculer $P(X>k)$.
\item Montrer propriété sans mémoire.
\end{enumerate}

%--------------------------------------
\textbf{Exercice 34 [PB ⋆⋆⋆⋆⋆]}  
On considère $X \sim \mathcal{P}(\lambda)$.

\begin{enumerate}
\item Vérifier que c’est une loi.
\item Calculer $E(X)$.
\item Calculer $Var(X)$.
\item Interpréter.
\end{enumerate}

%--------------------------------------
\textbf{Exercice 35 [PB ⋆⋆⋆⋆⋆]}  
Problème complet :

On observe un système de file d’attente modélisé par une loi de Poisson.

\begin{enumerate}
\item Modéliser le nombre d’arrivées.
\item Calculer probabilités.
\item Étudier l’espérance.
\item Étudier approximation normale.
\item Interpréter physiquement.
\end{enumerate}

%====================================================
\newpage
\section{Corrigé détaillé}

%====================================================
\subsection*{Exercice 1}

On lance un dé équilibré.

\textbf{1) Univers}

L’univers est l’ensemble des issues possibles :
\[
\Omega = \{1,2,3,4,5,6\}
\]

Chaque issue est équiprobable, donc :
\[
P(\{k\}) = \frac{1}{6}
\]

\textbf{2) Probabilité d’un nombre pair}

Les nombres pairs sont :
\[
A = \{2,4,6\}
\]

Donc :
\[
P(A) = \frac{3}{6} = \frac{1}{2}
\]

\textbf{3) Probabilité d’un nombre ≥ 4}

\[
B = \{4,5,6\}
\]

\[
P(B) = \frac{3}{6} = \frac{1}{2}
\]

\textbf{4) Indépendance}

On calcule :
\[
A \cap B = \{4,6\}
\]

\[
P(A \cap B) = \frac{2}{6} = \frac{1}{3}
\]

Or :
\[
P(A)P(B) = \frac{1}{2} \times \frac{1}{2} = \frac{1}{4}
\]

Donc :
\[
P(A \cap B) \neq P(A)P(B)
\]

Conclusion : les événements ne sont pas indépendants.

%====================================================
\subsection*{Exercice 2}

Soit $X \sim \mathcal{B}(1,p)$.

\textbf{1) Loi}

\[
P(X=1)=p,\quad P(X=0)=1-p
\]

\textbf{2) Espérance}

Par définition :
\[
E(X)=0 \cdot (1-p) + 1 \cdot p = p
\]

\textbf{3) Variance}

On calcule :
\[
E(X^2)=0^2(1-p)+1^2 p = p
\]

Donc :
\[
Var(X)=E(X^2)-E(X)^2 = p - p^2 = p(1-p)
\]

%====================================================
\subsection*{Exercice 3}

On répète 10 expériences de Bernoulli de paramètre $p=0.4$.

\textbf{1) Loi}

La variable $X$ suit une loi binomiale :
\[
X \sim \mathcal{B}(10,0.4)
\]

\textbf{2) Probabilité $X=4$}

\[
P(X=4)=\binom{10}{4}(0.4)^4(0.6)^6
\]

\textbf{3) Probabilité $X \geq 1$}

On utilise le complémentaire :
\[
P(X \geq 1)=1-P(X=0)
\]

\[
P(X=0)=(0.6)^{10}
\]

Donc :
\[
P(X \geq 1)=1-(0.6)^{10}
\]

\textbf{4) Espérance et variance}

\[
E(X)=np=10 \times 0.4 = 4
\]

\[
Var(X)=np(1-p)=10 \times 0.4 \times 0.6 = 2.4
\]

%====================================================
\subsection*{Exercice 11}

Soit $X$ une Bernoulli.

\textbf{1) Espérance}

\[
E(X)=0(1-p)+1p = p
\]

\textbf{2) Calcul de $E(X^2)$}

\[
E(X^2)=0^2(1-p)+1^2 p = p
\]

\textbf{3) Variance}

\[
Var(X)=E(X^2)-E(X)^2 = p - p^2 = p(1-p)
\]

%====================================================
\subsection*{Exercice 12}

Soient $X_1,...,X_n$ indépendantes de loi Bernoulli$(p)$.

\textbf{1) Loi}

Chaque variable prend 0 ou 1. La somme compte les succès :

\[
X = \sum X_i
\]

Donc $X$ suit une loi binomiale.

\textbf{2) Espérance}

Par linéarité :
\[
E(X)=\sum E(X_i)=np
\]

\textbf{3) Variance}

Indépendance ⇒
\[
Var(X)=\sum Var(X_i)=np(1-p)
\]

%====================================================
\subsection*{Exercice 33}

Soit $X$ géométrique.

\textbf{1) Probabilité}

\[
P(X=k)=(1-p)^{k-1}p
\]

\textbf{2) Calcul de $P(X>k)$}

\[
P(X>k)=\sum_{n=k+1}^{\infty}(1-p)^{n-1}p
\]

C’est une série géométrique :

\[
P(X>k)=(1-p)^k
\]

\textbf{3) Sans mémoire}

\[
P(X>k+n|X>k)=\frac{P(X>k+n)}{P(X>k)}=\frac{(1-p)^{k+n}}{(1-p)^k}=(1-p)^n
\]

Donc propriété sans mémoire.

%====================================================
\subsection*{Exercice 34}

Soit $X \sim \mathcal{P}(\lambda)$.

\textbf{1) Vérification loi}

\[
\sum_{k=0}^\infty \frac{\lambda^k e^{-\lambda}}{k!}
= e^{-\lambda} e^\lambda = 1
\]

\textbf{2) Espérance}

\[
E(X)=\sum k \frac{\lambda^k e^{-\lambda}}{k!}
\]

On utilise :
\[
\sum k \frac{\lambda^k}{k!} = \lambda e^\lambda
\]

Donc :
\[
E(X)=\lambda
\]

\textbf{3) Variance}

\[
Var(X)=\lambda
\]


\end{document}